\section{数据角色、基本假设与符号}
原始附件只读，以可见题面及可见数据说明解释字段。A1--A3 的综合质量是工程代理而非人工质量真值；A4+A5 用于配比拟合，A6--A11 用于已观察的同项目检验，A12--A15 属估算外推压力。B1 是 $N,D$ 主干拟合来源，B7 为半合成质量扩展，B2/B3 是形状检验，B4/B5 只比较来源内方向，B10 是 estimated。C2 的得分和 C4 的资源元数据不能仅凭相似模型名称视作同版本配对；C3 混合历史指标需分来源，C6 Loss 需逐来源处理。

令 $N$ 为十亿参数、$D$ 为十亿训练 token、$\boldsymbol p$ 为 17 维非负且和为 1 的训练域比例，$Q_A$ 为 A 侧质量代理，$Q_B$ 为 B 侧原生质量指标，$L$ 为来源指定验证 Loss，$C$ 与预算 $B$ 均为 FLOPs，$L_c$ 为上下文 token 数，$S$ 为六项 Benchmark 的等权百分制均分。A/B 的质量和 Loss 坐标是否相同并无观测证据；主跨题计算明示使用单调质量代理及相对配比乘子，并在敏感性分析中改变其幅度。任何参数、预测或配方只在相应支持域与这些条件下解释。
